\documentclass[12pt]{article}

%% Basic document formatting
\usepackage{amsmath, amsthm, amssymb, amsfonts}
\usepackage{mathtools}
\usepackage{xspace}
\usepackage{thmtools}
\usepackage{graphicx}
\usepackage{setspace}
\usepackage{fancyhdr}
\usepackage{titling}
\usepackage[left=0.4in,right=0.4in,top=1in,bottom=1in]{geometry}
\usepackage{float}
\usepackage{hyperref}
\usepackage{mathrsfs}
\usepackage{tabularx}
\usepackage[utf8]{inputenc}
\usepackage[english]{babel}
\usepackage{framed}
\usepackage[dvipsnames]{xcolor}
\usepackage{environ}
\usepackage{tcolorbox}
\tcbuselibrary{theorems,skins,breakable}

\usepackage{enumitem}
\setlist[enumerate]{leftmargin=*}
\setlist[enumerate,1]{labelindent=\parindent}
\setlist[enumerate,2]{labelindent=0pt}

% Blackboard Bold
\newcommand{\Z}{\mathbb{Z}}                 % integers
\newcommand{\Zpl}{\mathbb{Z}_{+}}           % positive integers
\newcommand{\Znn}{\mathbb{Z}_{\geq 0}}      % nonnegative integers. 
\newcommand{\Q}{\mathbb{Q}}                 % rationals
\newcommand{\Qpl}{\mathbb{Q}_{+}}           % positive Rationals
\newcommand{\R}{\mathbb{R}}                 % reals
\newcommand{\Rpl}{\mathbb{R}_{+}}           % positive Reals
\newcommand{\C}{\mathbb{C}}                 % complex numbers
\newcommand{\F}{\mathbb{F}}                 % field

% Words
\newcommand{\st}{\text{ such that }}
\newcommand{\wrt}{\text{ with respect to }}
\newcommand{\with}{\text{ with }}
\newcommand{\ie}{\text{, i.e. }}
\newcommand*{\ditto}{\text{---}\texttt{"}\text{---}}

% Operators
\newcommand{\abs}[1]{\left|#1\right|}                   % absolute value
\newcommand{\norm}[1]{\abs{\abs{#1}}}
\newcommand{\floor}[1]{\left\lfloor #1 \right\rfloor}   % floor
\newcommand{\ceil}[1]{\left\lceil #1 \right\rceil}      % ceiling
\newcommand{\powset}[1]{\mathscr{P}(#1)}

% Category Theory 
\renewcommand{\hom}{\text{Hom}}
\newcommand{\homspace}[3]{\hom_{#1}(#2, #3)}
\newcommand{\coproduct}[9]{
  \begin{tikzcd}[ampersand replacement=\&]
      \& \& #1 \arrow[dll, bend right, "#6"] \arrow[dl, "#5"] \\
   #4 \& #3 \arrow[l, "#9"] \\
      \& \& #2 \arrow[ull, bend left, "#8"] \arrow[ul, "#7"]
  \end{tikzcd}
}

\newcommand{\product}[9]{ 
  \begin{tikzcd}[ampersand replacement=\&]
      \& \& #3 \\
      #1 \arrow[r, "#7"] \arrow[urr, bend left, "#5"] \arrow[drr, bend right, "#6"] \& #2 \arrow[ur, "#8"] \arrow[dr, "#9"] \\ 
      \& \& #4
  \end{tikzcd}
}


% Algebra 
\newcommand{\Syl}[2]{\text{Syl}_{#1}{(#2)}}           % set of Sylow-p subgroups 
\newcommand{\<}{\langle}                % \<x\>, subgroup generated by x 
\renewcommand{\>}{\rangle}
\renewcommand{\index}[2]{[#1 : #2]}
\newcommand{\id}{\text{id}}             % identity element
\newcommand{\lhdneq}{%
  \mathrel{\ooalign{$\lneq$\cr\raise.22ex\hbox{$\lhd$}\cr}}}
\newcommand{\order}[1]{\text{o}(#1)}    % order of an element
\newcommand{\spec}{\operatorname{Spec}}
\newcommand{\rad}{\operatorname{rad}}   % radical of an ideal 
\newcommand{\sym}[1]{\mathfrak{S}_{#1}}
\newcommand{\aut}[1]{\text{Aut}(#1)} 
\newcommand{\gal}[2]{\text{Gal}(#1 / #2)} 
\newcommand{\inn}[1]{\text{Inn}(#1)} 
\let\oldcong\cong
\let\oldequiv\equiv
\renewcommand{\cong}{\oldequiv}
\renewcommand{\equiv}{\oldcong}

\newcommand{\triangleleftneq}{\mathrel{\ooalign{%
  $\trianglelefteq$\cr
  \hidewidth\raise0.06ex\hbox{$\not\mkern4mu$}\hidewidth\cr
}}}

\makeatletter % cycle
\newcommand{\cyc}[1]{(\mathbf{\cyc@process#1\relax})}
\def\cyc@process#1#2\relax{%
	#1%
	\ifx\relax#2\relax
	\else
		\,\cyc@process#2\relax
	\fi
}
\makeatother

\newcommand{\GL}[2]{\text{GL}_{#1}(#2)}                             % general linear group
\newcommand{\SL}[2]{\text{SL}_{#1}(#2)}                             % special linear group
\newcommand{\gl}[2]{\mathfrak{gl}_{#1}(#2)}
\renewcommand{\sl}[2]{\mathfrak{sl}_{#1}(#2)}
\newcommand{\so}[2]{\mathfrak{so}_{#1}(#2)}
\newcommand{\su}[1]{\mathfrak{su}(#1)}

% Linear Algebra
\newcommand{\bmat}[1]{\begin{bmatrix}#1\end{bmatrix}}   % bracketed matrix
\newcommand{\rank}{\operatorname{rank}}                 % rank
\newcommand{\nullity}{\operatorname{nullity}}           % nullity
\newcommand{\tr}{\operatorname{tr}}

% Combinatorics
\newcommand{\qnum}[1]{\left[#1\right]_q}                % q-analog
\newcommand{\qfact}[1]{\left[#1\right]!_q}              % q-analog factorial 
\newcommand{\qbinom}[2]{\binom{#1}{#2}_q}               % Gaussian binomial coefficient

% Topology/Analysis
\newcommand{\ball}[2]{\text{B}_{#1}(#2)}  % B_r(x): open r-balls around x
\newcommand{\diam}{\text{diam}}           % diamter of a set in metric space 
\newcommand{\lowint}[2]{\underline{\int_{#1}^{#2}}}
\newcommand{\upint}[2]{\overline{\int}_{#1}^{#2}}

% Misc Notation
\newcommand{\oneton}[1]{\{1, \ldots, #1\}}
\newcommand{\defeq}{\vcentcolon=}   % :=
\newcommand{\eqdef}{=\vcentcolon}   % =: 
\renewcommand{\bf}[1]{\textbf{#1}}
\newcommand{\im}{\operatorname{im}}
\newcommand{\fnrest}[2]{\left.#1\right|_{#2}}
\newcommand{\isom}{\overset{\sim}{\rightarrow}} 


\renewcommand{\thesection}{\arabic{section}.}
\renewcommand{\thesubsection}{(\alph{subsection})}

\newtheorem{claim}{Claim}
\newtheorem{theorem}{Theorem}
\newtheorem{proposition}{Proposition}
\newtheorem*{lemma}{Lemma}

%% Headers & title setup
\newcommand{\course}{Math140C}
\newcommand{\myname}{Jay Ser}
\setlength{\headheight}{14.5pt}
\pagestyle{fancy}
\fancyhf{}
\renewcommand{\headrulewidth}{0.4pt}
\lhead{\course}
\rhead{\myname}
\cfoot{\thepage}
\setlength{\droptitle}{-4em} 
\title{\course\ - HW \#3}
\author{\myname}
\date{2026.05.03} 

\begin{document}
\maketitle
\thispagestyle{fancy}

%------------------------------------------------------------------------------%

\newcommand{\Djf}[3]{[D_{#1}#2(#3)]}
\newcommand{\grad}[1]{[\nabla #1]}

\section*{Q23} 
$f(0, 1, -1) = 0 + 1 - 1 = 0$. 
Next, 
$$[f'(x, y_1, y_2)] = [2xy + e^x, x^2, 1 ],$$
so $\Djf{1}{f}{0, 1, -1} = 1$.
Namely, in the language of the implicit function theorem, $A_x$ is invertible. 
Hence the implicit function theorem applies; 
there is a neighborhood $W \subset \R^2$ of $(1, -1)$ and function $g: W \rightarrow \R$ such that $g(1, -1) = 0$ and  
$$f(g(y_1, y_2), y_1, y_2) = 0, \quad ((y_1, y_2) \in W)$$
Furthermore, the implicit function theorem gives the formula 
\begin{align*}
  [g'(1, -1)] & = -\Djf{1}{f}{0, 1, -1}^{-1} [\Djf{2}{f}{0, 1, -1}, \Djf{3}{f}{0, 1, -1}]  \\ 
              & = -1[0, 1] = [0, -1]
\end{align*}
In other words, $\Djf{1}{g}{1, -1} = 0$ and $\Djf{2}{g}{1, -1} = -1$. 

\section*{Q27}
\subsection{} 
\begin{align*}
  \limsup_{(x, y) \rightarrow (0, 0)} \abs{\frac{xy(x^2 - y^2)}{x^2 + y^2}}
  & \leq \lim_{(x, y) \rightarrow (0, 0)} \abs{\frac{x^3 y}{x^2 + y^2}} + \lim_{(x, y) \rightarrow (0, 0)} \abs{\frac{xy^3}{x^2 + y^2}} \\ 
  & = \lim_{(x, y) \rightarrow (0, 0)} \abs{xy}\abs{\frac{x^2}{x^2 + y^2}} + \lim_{(x, y) \rightarrow (0, 0)} \abs{xy}\abs{\frac{y^2}{x^2 + y^2}} \\ 
  & \leq 2 \lim_{(x,y ) \rightarrow (0, 0)} \abs{xy} = 0,
\end{align*}
So $\lim_{(x, y) \rightarrow (0, 0)} f(x, y) = 0$, meaning $f(x, y)$ is continuous at $(0, 0)$. 
Of course, $f$ is continuous away from $(0, 0)$ because it is rational.  
Away from $(0, 0)$, one can calculate using the quotient rule
\begin{align}
  \Djf{1}{f}{x, y} & = \frac{y(x^4 + 4x^2 y^2 -y^4)}{(x^2 + y^2)^2} \label{eq:27_D1f} \\ 
  \Djf{2}{f}{x, y} & = \frac{x^5 - 4x^3 y^2 - xy^4}{(x^2 + y^2)^2} \label{eq:27_D2f}
\end{align}
Also, 
\begin{align*}
  \lim_{h \rightarrow 0} \frac{f(h, 0) - f(0, 0)}{h} & = 0 \\  
  \lim_{h \rightarrow 0} \frac{f(0, h) - f(0, 0)}{h} & = 0
\end{align*}
shows $\Djf{1}{f}{0, 0} = 0$ and $\Djf{2}{f}{0, 0} = 0$. 
\begin{align*}
\limsup_{(x, y) \rightarrow (0, 0)} \abs{\frac{x^5 - 4x^3 y^2 - xy^4}{(x^2 + y^2)^2}}
  & \leq \lim_{(x, y) \rightarrow (0, 0)} \abs{x}\abs{\frac{x^4}{(x^2 + y^2)^2}} + \lim_{(x, y) \rightarrow (0, 0)} \abs{4x}\abs{\frac{x^2 y^2}{x^2 + y^2)^2}} \\ 
  &  \quad + \lim_{(x, y) \rightarrow (0, 0)} \abs{x}\abs{\frac{x^4}{(x^2 + y^2)^2}} \\ 
  & \leq \lim_{(x, y) \rightarrow (0, 0)} \abs{x} + \lim_{(x, y) \rightarrow (0, 0)} \abs{4x} + \lim_{(x, y) \rightarrow (0, 0)} \abs{x} = 0
\end{align*}
shows $D_{2}f$ is continuous at $(0, 0)$, and similar bounding shows $D_{1}f$ is continuous at $(0, 0)$ as well. 

\subsection{} 
\eqref{eq:27_D1f} and \eqref{eq:27_D2f} show that away from $(0, 0)$, $D_{12}f$ and $D_{21}f$ exist and are continuous. 
At $(0, 0)$,
\begin{align*}
  \Djf{21}{f}{0, 0} & = \lim_{h \rightarrow 0} \frac{\frac{h(-h^4)}{h^4} - 0}{h} = -1 \\ 
  \Djf{12}{f}{0, 0} & = \lim_{h \rightarrow 0} \frac{\frac{h^5}{h^4} - 0}{h} = 1
\end{align*}
Since $\Djf{21}{f}{0, 0} \neq \Djf{12}{f}{0, 0}$, one concludes that neither $D_{21}f$ nor $D_{12}f$ are continuous at $(0, 0)$. 

\subsection{} 
See (b). 

\section*{Q28}
For a fixed $t > 0$, $\varphi(x, t)$ is continuous on each piece. 
Further, the values of each piece match at $x = \sqrt{t}$ and $x = 2\sqrt{t}$, so $\varphi$ is continuous for all $x$ for this fixed $t > 0$. 
By symmetry, $\varphi$ is continuous at all $(x, t)$ for $t < 0$. 
For points at $t = 0$, notice that for every $(x, t)$ such that $t \neq 0$, 
$$\abs{\varphi(x, t)} \leq \sqrt{\abs{t}}$$
Hence in arbitrary small neighborhood around any point $(x, 0)$, the absolute value of $\varphi$ is arbitarily small. 
This shows that $\varphi$ is continuous on all of $\R^2$. 

Next, consider 
\begin{equation} \label{eq:28_partial}
  \lim_{h \rightarrow 0} \frac{\varphi(x, h) - 0}{h}
\end{equation}
By definition of $\varphi$ it is clear \eqref{eq:28_partial} is 0 when $x \leq 0$. 
When $x > 0$, it is also clear that as $h$ gets arbitrarily small, $x$ lands on the ``otherwise" part of $\varphi$, hence \eqref{eq:28_partial} is 0 once again. 
This shows $\Djf{2}{\varphi}{x, 0} = 0$. 

Assume $\abs{t} < 1/4$. 
Then $2\sqrt{t} < 1$, meaning $\varphi(x, t) = 0$ when $x \notin [0, 2\sqrt{t}]$ portion of $[-1, 1]$. 
\begin{align*}
  f(t) & = \int_{0}^{\sqrt{t}}xdx + \int_{\sqrt{t}}^{2\sqrt{t}} (-x + 2\sqrt{t})dx \\ 
       & = \frac{1}{2}\left.x^2\right|_{0}^{\sqrt{t}} + \left[-\frac{x^2}{2} + 2\sqrt{t}x\right]_{\sqrt{t}}^{2\sqrt{t}} \\ 
       & = \frac{t}{2} + (-2t + 4t - (-\frac{t}{2} + 2t)) \\ 
       & = t
\end{align*}
So $f(t) = t$ for $t \in [-1/4, 1/4]$. 
This gives $f'(0) = 1$. 
However, $\Djf{2}{\varphi}{x, 0} = 0$ for every $x$, so 
$$\int_{-1}^{1} \Djf{2}{\varphi}{x, 0} dx = 0.$$
In conclusion, 
$$\int_{-1}^{1} \Djf{2}{\varphi}{x, 0}dx \neq f'(0)$$

\end{document}
